%% ma: L10-B12-0018
%% lop: 10
%% chuong: 5
%% bai: 12
%% dang: 10.12.3
%% loai: TL
%% muc: TH
%% dap_an: "-"
%% nguon: "chuyen-de-cac-so-dac-trung-cua-mau-so-lieu-khong-ghep-nhom-toan-10-chuong-trinh-moi.pdf (GV. Phan Nhật Linh) (P01), Bài 01 Bài tập tự luận 18"
%% kien_thuc: "Lớp 10 Bài 12 (số gần đúng, sai số tuyệt đối, độ chính xác, sai số tương đối, quy tròn số)"
%% trang_thai: chinh-thuc
%% ghi_chu: "Đáp án gốc ý c) là 2,718282 (quy tròn đến hàng phần triệu, sai số ≈ 1,7.10^-7 vẫn nhỏ hơn d); kho dùng đúng quy tắc trong phần lý thuyết của tài liệu, được 2,7182818. Tự duyệt theo yêu cầu giáo viên 22/10/2026."
\begin{ex}
Số $e$ có giá trị gần đúng với $12$ chữ số thập phân là $2{,}718281828459$.

a) Giả sử lấy $2{,}7$ làm giá trị gần đúng của $e$. Chứng minh sai số tuyệt đối không vượt quá $0{,}02$ và sai số tương đối không vượt quá $0{,}75\%$.

b) Hãy quy tròn $e$ đến hàng phần nghìn.

c) Tìm số gần đúng của $e$ với độ chính xác $0{,}0000002$.
\loigiai{a) $\Delta=|e-2{,}7|\approx0{,}0182818<0{,}02$; $\delta=\dfrac{\Delta}{2{,}7}\approx0{,}677\%<0{,}75\%$.

b) $e\approx2{,}718$.

c) Chữ số khác $0$ đầu tiên của $d=0{,}0000002$ ở hàng thứ bảy sau dấu phẩy, nên quy tròn $e$ đến hàng đó: $e\approx2{,}7182818$ (sai số khoảng $3\cdot10^{-8}<0{,}0000002$).}
\end{ex}
